% The Slide Definitions
%document
\documentclass[10pt]{beamer}
%theme
\usetheme{metropolis}
% packages
\usepackage{color}
\usepackage{listings}
\usepackage[ngerman]{babel}
\usepackage[utf8]{inputenc}
\usepackage{multicol}


% color definitions
\definecolor{mygreen}{rgb}{0,0.6,0}
\definecolor{mygray}{rgb}{0.5,0.5,0.5}
\definecolor{mymauve}{rgb}{0.58,0,0.82}

\lstset{
    backgroundcolor=\color{white},
    % choose the background color;
    % you must add \usepackage{color} or \usepackage{xcolor}
    basicstyle=\footnotesize\ttfamily,
    % the size of the fonts that are used for the code
    breakatwhitespace=false,
    % sets if automatic breaks should only happen at whitespace
    breaklines=true,                 % sets automatic line breaking
    captionpos=b,                    % sets the caption-position to bottom
    commentstyle=\color{mygreen},    % comment style
    % deletekeywords={...},
    % if you want to delete keywords from the given language
    extendedchars=true,
    % lets you use non-ASCII characters;
    % for 8-bits encodings only, does not work with UTF-8
    frame=single,                    % adds a frame around the code
    keepspaces=true,
    % keeps spaces in text,
    % useful for keeping indentation of code
    % (possibly needs columns=flexible)
    keywordstyle=\color{blue},       % keyword style
    % morekeywords={*,...},
    % if you want to add more keywords to the set
    numbers=left,
    % where to put the line-numbers; possible values are (none, left, right)
    numbersep=5pt,
    % how far the line-numbers are from the code
    numberstyle=\tiny\color{mygray},
    % the style that is used for the line-numbers
    rulecolor=\color{black},
    % if not set, the frame-color may be changed on line-breaks
    % within not-black text (e.g. comments (green here))
    stepnumber=1,
    % the step between two line-numbers.
    % If it's 1, each line will be numbered
    stringstyle=\color{mymauve},     % string literal style
    tabsize=4,                       % sets default tabsize to 4 spaces
    % show the filename of files included with \lstinputlisting;
    % also try caption instead of title
    language = Python,
	showspaces = false,
	showtabs = false,
	showstringspaces = false,
	escapechar = ,
}

\def\ContinueLineNumber{\lstset{firstnumber=last}}
\def\StartLineAt#1{\lstset{firstnumber=#1}}
\let\numberLineAt\StartLineAt



\newcommand{\codeline}[1]{
	\alert{\texttt{#1}}
}


% Author and Course information
% This Document contains the information about this course.

% Authors of the slides
\author{Felix Döring, Felix Wittwer, Anton Obersteiner}

% Name of the Course
\institute{Python-Kurs}

% Fancy Logo 
\titlegraphic{\hfill\includegraphics[height=1.25cm]{../templates/fsr_logo_cropped}}



% Custom Bindings
% \newcommand{\codeline}[1]{
%	\alert{\texttt{#1}}
%}


% Presentation title
\title{Grundlagen}
\date{\today}
\setlength{\columnsep}{.4cm}

\begin{document}

\maketitle

\begin{frame}{Gliederung}
	\begin{multicols}{2}
		\setbeamertemplate{section in toc}[sections numbered]
		\tableofcontents
	\end{multicols}
\end{frame}


\section{Wiederholung}
\subsection{Gemeinsamkeiten und Unterschiede}
\begin{frame}{Unterschiede und Gemeinsamkeiten}
    \begin{multicols}{2}
        \lstinputlisting{resources/02_grundlagen/say_hi_print.py}
        \onslide<2->\lstinputlisting{resources/02_grundlagen/say_hi_return.py}
    \end{multicols}
    \begin{tabular}{r|l|l}
        \onslide<2->Eigenschaft & \texttt{print}-Version & \texttt{return}-Version \\ \hline
        \onslide<3->{Strings?    & \texttt{f"\{var\}"} Format-String & {\textcolor{red}{\texttt{+}}}-Operator } \\
        \onslide<4->{Argumente   & \texttt{greeting} hat Standard- & Beide Argumente sind \\
                    & Argument \texttt{"Hi"} & notwendig } \\
        \onslide<5->{Fehler?    & keine & Missing Argument } \\
        \onslide<6->{Effekt      & Funktion gibt aus & Funktion gibt zurück \\
                    &   & $\to$ in Datei geschrieben } \\
        \onslide<7->{Bedingung   & Nur wenn Skript & wird immer gemacht \\
                    & direkt ausgeführt wird & \\
                    & (\texttt{\_\_name\_\_ == "\_\_main\_\_"}) & } \\
    \end{tabular}
\end{frame}
\begin{frame}{Boilerplate}
	\lstinputlisting{resources/02_grundlagen/boilerplate.py}
	Ausführbare Python-Datei vs. Bibliothek $\to$ \alert{\tt main}-Funktion.\\
	Boilerplate unterscheidet: Main oder Bibliothek? \\
\end{frame}
\subsection{Fehler}
\begin{frame}{Auftreten von Fehlern}
    \begin{itemize}
        \item Wann bemerkt man Fehler?
        \item<2-> Syntax-Fehler am Anfang
        \item<3-> Alle anderen bei der Ausführung!
        \item<4-> Deswegen komplizierten Code Testen! (später)
    \end{itemize}
\end{frame}

\section{Scopes}
\subsection{Das Scoping Problem}
\begin{frame}{Das Problem mit dem Scope}
	\lstinputlisting{resources/02_grundlagen/scope.py}
    \onslide<2->{
	\textbf{\alert{Das Problem:}} \\[.25cm]
	Variablen sind zwar nach innen sichtbar, werden aber innerhalb der Funktion neu angelegt.}
\end{frame}

\begin{frame}{\tt global}
	\lstinputlisting{resources/02_grundlagen/scope_fix.py}
    Das ist meist Gaffa, keine dauerhafte Lösung! \\
    $\Rightarrow$ Variablen als Argumente übergeben
\end{frame}
\begin{frame}{Parameter?}
	\lstinputlisting{resources/02_grundlagen/scope_parameter.py}
    \glqq passing by value\grqq{} statt \glqq by reference\grqq{} :( \\
    zurückgeben $\Rightarrow$ Funktionaler Ansatz
\end{frame}
\begin{frame}{Funktionaler Ansatz}
	\lstinputlisting{resources/02_grundlagen/scope_functional.py}
    Alles übergeben und zurückgeben... Anstrengend \\
    sinnvoll strukturieren $\Rightarrow$ Objekte
\end{frame}

\section{Objekte}
\subsection{Klassen, Objekte und Attribute}
\begin{frame}{Klassen und Attribute}
	\begin{itemize}
		\item<1-> Klassen sind wie Schablonen
        \item<2-> Objekte sind konkrete Instanzen
	\end{itemize}
    \onslide<3->\lstinputlisting{resources/02_grundlagen/class.py}
    \onslide<3->{auf Deutsch: selbst festgelegt}
    %%don't do -> not shown
	% \begin{itemize}
	% 	\item Man kann Objekten jederzeit neue Attribute hinzufügen (auch \texttt{type}-Objekten)
	% 	\item Ist außerhalb des Initialisierers nicht empfehlenswert
	% \end{itemize}
\end{frame}

\begin{frame}{Attribute}
	Attribute setzen: meist im \textbf{Initialisierer} \alert{\tt \_\_init\_\_}
	\lstinputlisting{resources/02_grundlagen/attributes.py}
    auf Deutsch: selbst festgelegt
\end{frame}

\subsection{Methoden}
\begin{frame}{Methoden}
	\begin{itemize}
		\item Funktionen, die Teil von Klassen/Objekten sind
		\item<2-> erstes Argument \texttt{self} wird übergeben
        \item<3-> \texttt{    def method(self, arg2): ...}
        \item<3-> \texttt{obj.method(arg2)}
	\end{itemize}
\end{frame}

\subsubsection{Übung: Vektor}
\begin{frame}{Übung: Vektor-Klasse}
    \lstinputlisting{resources/02_grundlagen/vector_question.py}
\end{frame}
\begin{frame}{Übung: Vektor-Klasse}
    \lstinputlisting{resources/02_grundlagen/vector.py}
\end{frame}

\subsection{Super- und Subklassen}
\begin{frame}{Super- und Subklassen}
	Erben / Erweitern: mit \alert{\tt class subclass(superclass):} definieren
	\begin{itemize}
		\item neue Variablen und Methoden hinzufügen, auch alte überschreiben
		\item Superklasse mit \alert{\tt super()} aufrufen
    \end{itemize}
	\lstinputlisting{resources/02_grundlagen/inheritance.py}
\end{frame}
\subsubsection{Übung: Planet}
\begin{frame}{Übung: Planet-Klasse}
    \lstinputlisting{resources/02_grundlagen/planet_question.py}
\end{frame}
\begin{frame}{Übung: Planet-Klasse}
    \lstinputlisting{resources/02_grundlagen/planet.py}
\end{frame}
\begin{frame}{Übung: Planet-Klasse (Fortsetzung)}
    \lstinputlisting{resources/02_grundlagen/planet_2.py}
\end{frame}

\subsection{Spezielle Methoden}
\begin{frame}{Spezielle Methoden}
	von Python intern verwendet

	\begin{description}
        \item<1->[\tt\_\_init\_\_] Bei Instanziierung aufgerufen
		\item<2->[\tt\_\_del\_\_] Bei Löschung aufgerufen (selten)
        \item<3->[\tt\_\_str\_\_] Für Ausgabe als String
        \item<4->[\tt\_\_add\_\_] Funktion für \alert{\tt +}-Operator
        \item<5->[\tt\_\_iadd\_\_] Funktion für \alert{\tt +=}-Operator (muss self zurückgeben)
	\end{description}
\end{frame}
\begin{frame}{Beispiel: Mein Vektor}
    \lstinputlisting{resources/02_grundlagen/Vector.py}
\end{frame}


\end{document}
